\chapter{Forensic Medicine \& Medical Ethics}

\medterm{Adipocere (Grave Wax)}-- Waxy, grayish-white substance formed postmortem in warm, moist, anaerobic environments (water, damp soil) by the hydrolysis and hydrogenation of body fat into free fatty acids; retards putrefaction and preserves anatomical details for months to years.

\medterm{Algor Mortis}-- Post-mortem cooling of the body following death until it reaches ambient environmental temperature. Serves as one of the physical estimators of the post-mortem interval (time since death). Under standard conditions ($20^\circ\text{C}$ ambient temperature), the body cools at a rate of approximately $0.75\text{--}1.0^\circ\text{C}$ ($1.5^\circ\text{F}$) per hour during the first 12 hours, and ~ $0.5^\circ\text{C}$ per hour thereafter. Influenced by ambient temperature, body habitus (obesity slows cooling), clothing, humidity, air movement, and immersion in water. Rectal temperature is recorded using a chemical thermometer.

\medterm{Antemortem Injury}-- Injury inflicted upon a body prior to death, characterized microscopically and macroscopically by physiological tissue responses including hemorrhage, blood clotting, inflammatory cell infiltration, tissue edema, and active healing reactions (differentiating it from postmortem artifact).

\medterm{Autopsy (Medicolegal)}-- Systematic external examination and surgical dissection of a deceased body performed under legal authority by a forensic pathologist to determine the Cause, Manner, and Mechanism of Death, document injuries, collect evidence, and assist judicial investigations.

\medterm{Autopsy (Post-Mortem Examination)}-- Specialized surgical procedure performed on a deceased body to determine the cause of death, mechanism of death, manner of death (natural, accidental, suicide, homicide, or undetermined), state of health prior to death, and to document injuries or disease processes. Types: Clinical Autopsy (hospital-based, requiring family consent, to study disease processes and audit medical care) and Forensic (Medico-Legal) Autopsy (ordered by a coroner, medical examiner, or magistrate under statutory authority, without requiring family consent, in cases of sudden, violent, suspicious, or unexplained death). Procedure involves external examination, full internal dissection of cranial, thoracic, and abdominal cavities, histopathology, toxicological sampling, and formal report generation.

\medterm{Beneficence and Non-Maleficence}-- Core biomedical ethical principles guiding clinical practice and medical research. Beneficence requires healthcare providers to act in the best interest of the patient, taking positive steps to promote health, prevent harm, and relieve suffering. Non-Maleficence (expressed by the Latin phrase \textit{primum non nocere}—"first, do no harm") obligates practitioners to avoid inflicting intentional harm, unnecessary pain, or undue risk on patients. Ethical dilemmas arise when balancing the potential benefits of a treatment against its inherent risks and side effects.

\medterm{Bioethics}-- Interdisciplinary study of the ethical, legal, philosophical, and social implications of advancements in medicine, biology, healthcare technology, and biotechnology. Key topics: autonomy, informed consent, end-of-life decision-making (euthanasia, advance directives, do-not-resuscitate orders), reproductive technologies (surrogacy, gene editing, IVF), organ transplantation ethics, allocation of scarce healthcare resources, genetic privacy, and clinical trial ethics.

\medterm{Cadaveric Spasm}-- Rare form of instantaneous post-mortem stiffening of a specific group of voluntary muscles at the exact moment of death, without an intervening period of muscle relaxation. Unlike normal rigor mortis (which develops 2--4 hours post-mortem), cadaveric spasm occurs immediately, trapping objects held in the hand at death (e.g., a weapon in suicide, grass/twigs held by a drowning victim). Pathophysiology: severe physical or emotional stress combined with extreme exhaustion immediately preceding death, leading to rapid ATP depletion. Highly significant in forensic investigations for establishing suicide or struggle.

\medterm{Capacity and Competence}-- Decision-making capacity is a clinical evaluation performed by a physician to determine whether a patient understands the nature, consequences, risks, and alternatives of a specific proposed medical decision at a given time. Competence is a formal legal determination made by a judge regarding an individual's overall legal ability to manage their affairs or make decisions. A patient with capacity has the legal and moral right to grant or refuse consent for medical procedures, even if that refusal leads to harm or death.

\medterm{Carbon Monoxide Poisoning}-- Toxicology emergency resulting from inhalation of CO gas, which binds hemoglobin with $200\times$ higher affinity than oxygen to form Carboxyhemoglobin (COHb), shifting oxygen-hemoglobin dissociation curve left. Classical autopsy finding: cherry-red discoloration of skin, blood, mucosal surfaces, and livor mortis.

\medterm{Cause of Death}-- The specific disease, injury, or toxic agent responsible for initiating the lethal physiological sequence leading directly to death (e.g., Gunshot wound to the head, Acute myocardial infarction).

\medterm{Chain of Custody}-- Meticulous paper trail and chronological documentation tracking the seizure, custody, control, transfer, analysis, and disposition of physical and electronic evidence in medico-legal investigations (e.g., blood samples for toxicology, bullets, sexual assault kits, biological fluids). Essential in forensic science to demonstrate that evidence presented in a court of law is authentic, uncontaminated, and has remained unaltered since collection. Any breach in the chain of custody can render evidence inadmissible.

\medterm{Confidentiality}-- The ethical and legal duty of healthcare professionals to protect patient information from unauthorized disclosure. Founded on the principle of respect for patient autonomy and trust, confidentiality is a core requirement of medical ethics (GMC guidelines) and law (GDPR, Data Protection Act, Caldicott principles). Information should only be shared with the patient's explicit consent, or when required by law (public health reporting, court order), or when necessary to protect the patient or others from serious harm (safeguarding, GMC duty of candor). Breach of confidentiality can result in disciplinary action, professional sanctions, and legal liability. Exceptions: serious communicable disease reporting (HIV, TB), child protection concerns, fitness to drive notification, and terrorism prevention. The duty of confidentiality continues after death (access to records of the deceased requires executor or next-of-kin consent).

\medterm{Defense Mechanisms (Forensic)}-- Physical evidence patterns on a body demonstrating active resistance against an assailant prior to incapacitation.

\medterm{Defense Wounds}-- Injuries sustained by a victim attempting to ward off an attack, typically found on the flexor surfaces of hands, palms, wrists, and forearms (incised/stab wounds from knives or contusions/abrasions from blunt force).

\medterm{Defense Wounds (Inscribed)}-- Patterned incised or lacerated trauma localized to upper extremities indicating active physical defense.

\medterm{Diatom Test}-- Forensic technique used to support the diagnosis of drowning; microscopic silica-shelled diatoms present in drowning water enter lungs, penetrate alveolar capillaries, and circulate to distant closed organs (femur bone marrow, brain) only if active circulation was present at time of submersion.

\medterm{Drowning (Dry/Wet Signs)}-- Cause of death resulting from submersion or immersion in liquid causing primary respiratory impairment. Wet Drowning (85--90\%): fluid enters the lungs, disrupting surfactant, causing alveolar collapse, ventilation-perfusion mismatch, hypoxia, cardiac arrest, and characteristic voluminous frothy foam emerging from the mouth and nostrils (Paltauf's foam). Dry Drowning (10--15\%): severe laryngospasm triggered by initial water contact blocks fluid from entering the lungs, causing death by asphyxiation. Post-mortem signs: washerwoman skin (wrinkling of skin on palms and soles), cadaveric spasm holding debris, petechial hemorrhages in temporal bones (Wydler's sign), and water in the stomach and sphenoid sinus.

\medterm{Euthanasia}-- The deliberate act of ending a person's life to relieve suffering, typically performed by a physician. Active euthanasia involves directly administering a lethal substance (e.g., barbiturate overdose followed by neuromuscular blockade) at the patient's explicit request. Voluntary euthanasia is performed with the patient's informed consent. Physician-assisted suicide (PAS) involves the physician prescribing a lethal dose of medication that the patient self-administers. Euthanasia is legal in a growing number of countries and jurisdictions under strict regulatory conditions (Belgium, Netherlands, Luxembourg, Canada, Colombia, Spain, New Zealand, and several Australian states, as well as PAS in Switzerland, Germany, Austria, and several US states). It is distinct from palliative sedation (using medication to relieve refractory suffering at the end of life, not intended to hasten death) and withdrawal of life-sustaining treatment (allowing natural death).

\medterm{Euthanasia and Assisted Dying}-- Ethical and legal concepts regarding intentional termination of life to relieve intractable pain and suffering. Active Euthanasia: deliberate administration of a lethal agent (e.g., lethal dose of barbiturates) by a physician to terminate a patient's life upon their explicit request. Passive Euthanasia: withdrawal or withholding of life-sustaining treatment (e.g., mechanical ventilation, artificial nutrition), allowing natural death to occur. Physician-Assisted Suicide (Assisted Dying): a physician prescribes a lethal dose of medication that the patient self-administers. Legal status varies widely by international jurisdiction.

\medterm{Exhumation}-- Lawful disinterment of a buried corpse from a grave for medicolegal autopsy examination, re-examination, toxicological analysis, or paternity/DNA identification.

\medterm{Forensic Entomology}-- Application of insect biology (fly blowfly larva succession—Calliphoridae) to medicolegal investigations; analyzing insect species colonization and growth stages on a decomposing corpse provides accurate estimates of minimum Postmortem Interval (PMI).

\medterm{Forensic Medicine (Legal Medicine)}-- Branch of medicine that applies medical, biological, and scientific knowledge to legal questions and the administration of justice. Sub-fields: Forensic Pathology (determination of cause and manner of death via autopsy), Clinical Forensic Medicine (evaluation of living victims of sexual assault, domestic abuse, child abuse, torture, and blood alcohol/drug testing in drivers), Forensic Toxicology (detection of drugs and poisons in legal contexts), Forensic Psychiatry (evaluation of criminal responsibility and competency to stand trial), and Forensic Odontology (identification via dental records).

\medterm{Forensic Odontology}-- Branch of forensic science that applies dental knowledge to legal investigations. Primary applications: human identification of mass disaster victims, badly decomposed, charred, or skeletal remains by comparing ante-mortem dental records (X-rays, charts, fillings, prosthetics) with post-mortem dental findings; analysis of bite mark impressions on victims or food substances in violent crimes; and age estimation of living or deceased individuals based on dental eruption and tooth wear.

\medterm{Forensic Pathology}-- Medical subspecialty practiced by forensic pathologists (coroners/medical examiners) who investigate sudden, unexpected, violent, or suspicious deaths. Forensic pathologists perform autopsies, examine crime scenes, interpret wound patterns (gunshot wounds, stab wounds, blunt force trauma), estimate post-mortem intervals, determine cause and manner of death, and provide expert witness testimony in legal proceedings.

\medterm{Forensic Psychiatry}-- The interface of psychiatry and the legal system. It involves evaluations for criminal responsibility, competency to stand trial, malingering assessment, risk assessment for violence, and civil commitment. Forensic psychiatrists serve as expert witnesses and provide opinions on legal questions. Malingering is the intentional production of false or grossly exaggerated symptoms for external gain.

\medterm{Gettler Test}-- Historical quantitative chemical test for drowning, comparing chloride concentration in blood harvested from the left and right heart ventricles.

\medterm{Hanging and Strangulation Signs}-- Medico-legal findings resulting from mechanical constriction of the neck structures. Hanging: neck constriction by a ligature tightened by the weight of the victim's body. Ligature mark: pale, hard, parchment-like groove running obliquely upward toward the knot, usually situated above the thyroid cartilage. Strangulation: neck constriction by a ligature tightened by external force (ligature strangulation, horizontal mark around neck) or by hands/fingers (manual strangulation/throttling, showing fingernail abrasions and contusions). Internal autopsy findings: fracture of hyoid bone and thyroid cartilage cornua, intramuscular hemorrhages of neck muscles, mucosal petechiae, and Tardieu spots.

\medterm{Hanging (Complete vs Incomplete)}-- Form of violent asphyxial death produced by suspension of the body by a ligature around the neck. Complete Hanging: feet off ground, total body weight exerts pressure. Incomplete Hanging: part of body touches ground. Asphyxia results from occlusion of jugular veins ($2\text{ kg}$ force) or carotid arteries ($5\text{ kg}$ force).

\medterm{Hesitation Wounds}-- Multiple superficial, tentative parallel incised cuts surrounding a deeper primary fatal cut wound, classically observed in suicides involving sharp instruments (wrists, neck).

\medterm{Incised Wound}-- Clean-cut wound produced by a sharp-edged instrument (scalpel, knife, glass) that is longer on the skin surface than it is deep, featuring clean edges without tissue bridging or marginal abrasion.

\medterm{Informed Consent}-- Autonomous voluntary authorization given by a competent patient to undergo a specific medical intervention, procedure, or clinical trial. Core elements required for valid informed consent: (1) Disclosure of adequate information (nature of procedure, benefits, material risks, alternatives, and consequences of refusal), (2) Patient comprehension of the disclosed information, (3) Decision-making capacity of the patient, and (4) Voluntariness (freedom from coercion, manipulation, or undue influence). Exceptions: emergency situations where immediate intervention is required to save life or prevent severe harm and patient lacks capacity.

\medterm{Injury Classification (Abrasion, Laceration, Contusion, Incised Wounds)}-- Forensic categorization of mechanical wounds. Abrasion: superficial friction injury affecting only the epidermis (e.g., scrape, graze, pattern abrasion). Contusion (Bruise): extravasation of blood into subcutaneous tissues caused by blunt trauma crushing blood vessels while leaving overlying skin intact; changes color over time (red/purple $\rightarrow$ blue/black $\rightarrow$ green $\rightarrow$ yellow). Laceration: tear or split of skin and deep tissues caused by blunt impact or crushing force, characterized by irregular, bruised edges, tissue bridging (intact nerves/vessels across the gap), and hair bulbs crushed. Incised Wound (Cut): sharp force injury caused by a blade, longer than it is deep, with clean-cut, unbruised margins and no tissue bridging.

\medterm{Laceration (Forensic)}-- Tear or disruption of tissue produced by blunt force trauma exceeding skin elasticity, featuring irregular bruised edges, tissue bridging (intact nerves, vessels across wound floor), and hair crushing.

\medterm{Living Will (Advance Statements)}-- A legally binding written document that specifies a person's preferences for medical treatment if they become unable to make decisions for themselves in the future (loss of capacity). Living wills are a type of advance directive (or advance decision/advance statement in the UK). They allow individuals to refuse specific treatments (e.g., cardiopulmonary resuscitation, mechanical ventilation, artificial nutrition and hydration, antibiotics) in specified future scenarios (e.g., terminal illness, permanent unconsciousness/vegetative state, severe dementia). Legal requirements: the document must be signed and witnessed, signed in the presence of a witness, and it must clearly state the treatments being refused and the circumstances in which the refusal applies. In the UK: the Mental Capacity Act 2005 provides the legal framework for advance decisions. They are legally binding for refusing life-sustaining treatment if they are valid and applicable to the current circumstances. A valid advance decision can be overridden if the person has capacity and changes their mind when the situation arises. Related concepts: (1) Lasting Power of Attorney (LPA): allowed in the UK from 2007 for health and welfare decisions as well as decisions about property and financial affairs (in England and Wales), giving a named person legal authority to make decisions on behalf of the incapacitated patient. (2) Do Not Attempt Resuscitation (DNAR/DNACPR): a medical order signed by a senior clinician indicating that CPR should not be attempted in the event of cardiac arrest. (3) Physician Orders for Life-Sustaining Treatment (POLST): a medical order form used in the US that translates patient preferences into actionable medical orders. Advance care planning (ACP) discussions are encouraged to clarify patient values, goals, and preferences for future medical care.

\medterm{Livor Mortis (Lividity)}-- Post-mortem purplish-red discoloration of dependent body parts resulting from gravitational settling of blood in uncompressed capillaries and small veins following cessation of circulation. Onset: begins within 30 minutes to 2 hours after death; becomes fully developed and "fixed" (no longer blanches upon digital pressure due to hemolysis and extravasation) after 8--12 hours. Color variations provide forensic clues: cherry-red (carbon monoxide poisoning), bright pink (cyanide, cold exposure), brown (methemoglobinemia). Shifting of lividity patterns indicates that the body was moved post-mortem.

\medterm{Livor Mortis (Postmortem Lividity)}-- Purplish-red discoloration of dependent body parts following death, caused by gravitational pooling of blood in uncompressed capillaries and venules. Starts within 30 min--2 hours, becomes fixed after 6--12 hours. Color changes assist in diagnosing cause of death (cherry-red in CO, bright red in cold exposure/cyanide, chocolate brown in methemoglobinemia).

\medterm{Manner of Death}-- Legal classification of the circumstances surrounding a death, categorized into 5 official types: Natural, Accident, Suicide, Homicide, or Undetermined.

\medterm{Mechanism of Death}-- The specific physiological or biochemical derangement caused by the cause of death that leads to cessation of life (e.g., Exsanguination, Cardiac arrhythmia, Increased intracranial pressure).

\medterm{Medical Negligence (Malpractice)}-- Legal tort occurring when a healthcare provider fails to exercise the standard of reasonable care, skill, or diligence expected of a prudent practitioner of similar standing, resulting in injury, harm, or death to a patient. Legal elements required to establish medical negligence (the "4 Ds"): (1) Duty of care established between provider and patient, (2) Dereliction (breach) of that duty, (3) Direct causation (the breach directly caused the injury), and (4) Damages (measurable physical, financial, or emotional harm suffered by the patient).

\medterm{Mummification}-- Postmortem alteration occurring in hot, dry, arid, ventilated environments where rapid desiccation of tissues occurs before putrefaction can proceed, resulting in dry, leathery, dark brown skin clinging to underlying skeleton.

\medterm{Post-Mortem}-- An examination of a body after death to determine the cause of death, identify disease processes, and evaluate the effects of medical or surgical treatment. Also called autopsy, necropsy. Types: hospital (clinical) post-mortem — performed with consent to investigate the cause of death, assess treatment, or advance medical knowledge; medico-legal (coroner's/forensic) post-mortem — ordered by a coroner or legal authority when the death is sudden, unexpected, suspicious, or unnatural (required by law; consent not needed). Procedure: external examination, internal examination (thoracic, abdominal, and cranial cavities), histopathology (tissue sampling), toxicology (blood, urine, vitreous humor), microbiology, and other specialized tests. The report includes: cause of death (immediate, underlying, and contributing factors), and relevant pathological findings. Post-mortem examination is essential for quality assurance in medicine (verification of clinical diagnosis, identification of missed diagnoses, and audit of surgical and medical care). In obstetrics, perinatal post-mortem is crucial for understanding stillbirth, congenital anomalies, and maternal death.

\medterm{Postmortem Artifact}-- Any structural, chemical, or physical change introduced into a body after death (e.g., embalming artifacts, insect activity, postmortem animal scavenging, resuscitation trauma) that mimics antemortem disease or trauma.

\medterm{Postmortem Interval (PMI)}-- Estimated elapsed time between actual death and the time of autopsy or body discovery, calculated utilizing algor mortis, livor mortis, rigor mortis, vitreous potassium levels, and insect succession patterns.

\medterm{Punctate Hemorrhages (Tardieu Spots)}-- Small, pinpoint sub-ecchymotic hemorrhages (petechiae) resulting from the rupture of small capillaries under conditions of intense mechanical venous engorgement and asphyxia. Commonly observed on conjunctivae, sclerae, facial skin, epicardium, pleural surfaces, and thymus in deaths due to mechanical asphyxia (strangulation, hanging, suffocation, traumatic asphyxia). First described by Auguste Ambroise Tardieu in 1859.

\medterm{Putrefaction}-- Postmortem decomposition of body tissues brought about by bacterial and enzymatic action, characterized by green discoloration of lower abdomen (cecum), marbling of cutaneous veins, gas distension, foul odor, and liquefaction.

\medterm{Res Ipsa Loquitur}-- Latin legal doctrine meaning "the thing speaks for itself." Applied in medical malpractice cases where an injury is of a nature that ordinarily does not occur in the absence of negligence, the instrumentality causing harm was under the exclusive control of the defendant, and the plaintiff did not contribute to the injury. Examples: retaining a surgical sponge or instrument inside a patient's abdomen, or amputating the wrong limb. When invoked, it creates a rebuttable presumption of negligence, shifting the burden of proof to the defendant.

\medterm{Rigor Mortis}-- Post-mortem stiffening of body muscles resulting from loss of ATP synthesis following death. Without ATP, actin and myosin filaments remain permanently cross-linked in a rigid state. Onset: detectable in small muscles (jaw, eyelids) within 2--4 hours, becomes fully established in all voluntary muscles by 12 hours, remains fixed for 12 hours, and gradually resolves over the subsequent 12--24 hours as autolytic proteolytic enzymes break down muscle proteins ("rule of 12s": 12 hours onset, 12 hours duration, 12 hours resolution). Accelerated by high ambient temperature, fever, or violent exertion prior to death.

\medterm{Skeletonization}-- Final stage of postmortem decomposition where soft tissues, organs, and skin have completely decayed or been scavenged, leaving only clean osseous skeletal remains.

\medterm{Stippling (Tattooing)}-- Patterned dermal abrasions and punctate hemorrhages surrounding a firearm entrance wound, caused by unburned and partially burned gunpowder grains striking the skin at close range ($15\text{--}75\text{ cm}$).

\medterm{Sudden Infant Death Syndrome (SIDS)}-- Sudden, unexpected death of an infant under 1 year of age that remains unexplained after a thorough medico-legal investigation, including a complete autopsy, crime scene investigation, and review of clinical history. Peak incidence: 2--4 months of age. Pathophysiology: "triple-risk model" involving a vulnerable infant (brainstem respiratory control deficit), a critical developmental period, and an extrinsic stressor (prone sleeping position, soft bedding, exposure to cigarette smoke, overheating). Preventive campaign: "Back to Sleep" (placing infants supine for sleep) has reduced SIDS incidence by over 50\% globally.

\medterm{Thanatology}-- Scientific study of death, dying, post-mortem phenomena, and the psychological, social, and ethical aspects of bereavement. In forensic medicine, thanatology focuses on establishing clinical and legal definitions of death (brain death vs. somatic death), post-mortem changes (algor mortis, livor mortis, rigor mortis, decomposition, adipocere, mummification), and estimating the post-mortem interval (PMI).

